% ============================================================
\section{Các mô hình dự đoán}
\subsection{Hồi quy logistic}
\subsubsection{Khái niệm}
% ============================================================
 
Hồi quy logistic là một trong những thuật toán học máy có giám sát được sử dụng rộng rãi nhất trong lĩnh vực phân loại dữ liệu. Mặc dù mang tên hồi quy (\textit{regression}), mô hình này về bản chất là một phương pháp phân loại, được thiết kế để dự đoán xác suất thuộc về một lớp nhất định của biến mục tiêu rời rạc. Thuật toán này lần đầu tiên được giới thiệu trong lĩnh vực thống kê bởi nhà thống kê học David Cox vào năm 1958 và sau đó được tích hợp rộng rãi vào các hệ thống học máy hiện đại. \\
 
Về mặt toán học, hồi quy logistic là một mô hình tuyến tính tổng quát (\textit{Generalized Linear Model} - GLM), trong đó hàm liên kết là hàm logit. Mô hình này ánh xạ một tổ hợp tuyến tính của các biến đầu vào sang một giá trị xác suất nằm trong khoảng $(0,\,1)$, nhờ đó cho phép đưa ra quyết định phân lớp dựa trên ngưỡng xác suất được xác định trước. \\

Mô hình này được xây dựng nhằm giải quyết các bài toán phân loại nhị phân, tức là bài toán trong đó biến phụ thuộc chỉ nhận một trong hai giá trị, thường được ký hiệu là $0$ (âm tính) và $1$ (dương tính). Chẳng hạn, trong lĩnh vực y tế, mô hình có thể được sử dụng để dự đoán liệu một bệnh nhân có mắc bệnh tiểu đường hay không dựa trên các chỉ số sinh học. Dự đoán xác suất ghi bàn của một cú sút (vào hoặc không vào) trong lĩnh vực thể thao bóng đá. Hoặc trong tài chính, mô hình giúp xác định khả năng khách hàng có vỡ nợ hay không.\\

Ngoài bài toán nhị phân, hồi quy logistic có thể được mở rộng sang bài toán phân loại đa lớp, tức là bài toán trong đó biến mục tiêu $y$ có thể nhận một trong $K > 2$ giá trị rời rạc thay vì chỉ hai giá trị (0 hoặc 1) như thông thường. Bằng cách tiếp cận theo các hướng phổ biến như phân loại nhiều lớp với mỗi lớp là một bài toán nhị phân độc lập (One-vs-Rest) hoặc xây dựng các bộ phân loại trên từng cặp lớp (One-vs-One), hoặc mở rộng trực tiếp hàm sigmoid thành hàm softmax, hồi quy logistic có thể được áp dụng hiệu quả cho các bài toán phân loại đa lớp. Tuy nhiên, trong phạm vi luận văn này, trọng tâm được đặt vào biến thể nhị phân, dạng phổ biến và nền tảng nhất của mô hình.\\
% ============================================================

% ============================================================
\subsubsection{Hàm Sigmoid}
% ============================================================
 
Hàm sigmoid (còn gọi là hàm logistic) là trái tim của mô hình hồi quy logistic. Hàm này được định nghĩa như sau:
 
\begin{equation}
  \sigma(z) = \frac{1}{1 + e^{-z}}
  \label{eq:sigmoid}
\end{equation}
 
\noindent trong đó $z$ là giá trị đầu vào thực bất kỳ và $e$ là cơ số của logarithm tự nhiên.\\
 
Hàm sigmoid có một số tính chất toán học đặc biệt quan trọng. Thứ nhất, miền giá trị của hàm nằm hoàn toàn trong khoảng $(0,\,1)$, điều này khiến đầu ra của hàm có thể được diễn giải một cách tự nhiên như xác suất. Thứ hai, hàm có dạng chữ S đơn điệu tăng, nghĩa là khi $z \to +\infty$ thì $\sigma(z) \to 1$, và khi $z \to -\infty$ thì $\sigma(z) \to 0$. Thứ ba, tại $z = 0$, hàm cho giá trị $\sigma(0) = 0{,}5$, tương ứng với trạng thái không chắc chắn hoàn toàn giữa hai lớp. Ngoài ra, đạo hàm của hàm sigmoid có dạng thuận lợi cho tính toán:\\

\begin{equation}
  \sigma'(z) = \sigma(z)\bigl(1 - \sigma(z)\bigr)
  \label{eq:sigmoid_deriv}
\end{equation}
 
\noindent điều này giúp đơn giản hóa đáng kể quá trình tối ưu hóa tham số bằng phương pháp lặp cập nhật tham số theo hướng giảm nhanh nhất của hàm mất mát \textit{(gradient descent)} (được đề cập tại mục \ref{sec:gradient}).\\
 
Về mặt ý nghĩa thực tiễn, hàm sigmoid đóng vai trò là cầu nối chuyển đổi một đầu ra tuyến tính không bị giới hạn thành một xác suất có ý nghĩa thống kê, khắc phục nhược điểm chính của hồi quy tuyến tính khi áp dụng trực tiếp cho bài toán phân loại.\\

\subsubsection{Nguyên lý hoạt động}
% ============================================================
 
Hồi quy logistic hoạt động dựa trên ý tưởng cốt lõi là ước lượng xác suất có điều kiện $P(y = 1 \mid \mathbf{x})$, tức là xác suất để quan sát $\mathbf{x}$ thuộc về lớp dương tính. Thay vì dự đoán trực tiếp nhãn lớp, mô hình tính toán một giá trị xác suất liên tục, sau đó áp dụng một ngưỡng quyết định (\textit{threshold}) (thường là $0{,}5$) để gán nhãn cuối cùng.\\
 
Quá trình hoạt động của mô hình có thể được mô tả qua ba bước chính. Trước tiên, mô hình tính toán tổ hợp tuyến tính $z = \mathbf{w}^{T}\mathbf{x} + b$ từ các đặc trưng đầu vào. Tiếp theo, giá trị $z$ được đưa qua hàm sigmoid để chuyển đổi thành xác suất $\hat{p} = \sigma(z)$. Cuối cùng, dựa trên $\hat{p}$, mô hình đưa ra quyết định phân lớp theo ngưỡng đã định.\\
 
Mô hình hồi quy logistic được định nghĩa hoàn chỉnh thông qua công thức sau:
 
\begin{equation}
  \hat{p} = P(y = 1 \mid \mathbf{x})
           = \sigma\!\left(\mathbf{w}^{T}\mathbf{x} + b\right)
           = \frac{1}{1 + e^{-(\mathbf{w}^{T}\mathbf{x} + b)}}
  \label{eq:logreg}
\end{equation}
 
\noindent Trong đó:
 
\begin{itemize}
  \item $\mathbf{x} = (x_1, x_2, \ldots, x_n)^{T} \in \mathbb{R}^{n}$: Vector đặc trưng đầu vào gồm $n$ biến độc lập, mô tả các thuộc tính của mẫu dữ liệu cần phân loại.
 
  \item $\mathbf{w} = (w_1, w_2, \ldots, w_n)^{T} \in \mathbb{R}^{n}$: Vector trọng số - các tham số học được từ dữ liệu huấn luyện. Mỗi $w_i$ thể hiện mức độ ảnh hưởng của đặc trưng $x_i$ đến xác suất dự đoán. Trọng số dương đồng nghĩa với việc đặc trưng tương ứng làm tăng xác suất thuộc lớp $1$, trong khi trọng số âm thể hiện tác động ngược lại.
 
  \item $b \in \mathbb{R}$: Hệ số chệch - một hằng số cho phép mô hình linh hoạt trong việc dịch chuyển đường quyết định khỏi gốc tọa độ, giúp mô hình khớp tốt hơn với dữ liệu thực tế.
 
  \item $z = \mathbf{w}^{T}\mathbf{x} + b$: Giá trị tổng hợp tuyến tính, còn được gọi là \textit{log-odds} hay \textit{logit}, biểu thị logarithm của tỉ lệ xác suất:
  \[
    z = \log\!\left(\frac{\hat{p}}{1 - \hat{p}}\right)
  \]
 
  \item $\hat{p}$: Xác suất ước lượng rằng mẫu $\mathbf{x}$ thuộc về lớp dương tính ($y = 1$).
\end{itemize}
 
Sau khi tính được xác suất $\hat{p}$, mô hình đưa ra quyết định phân lớp dựa trên ngưỡng quyết định $\tau$ (thường là $0{,}5$) theo quy tắc:
 
\begin{equation}
  \hat{y} =
  \begin{cases}
    1 & \text{nếu } \hat{p} \geq \tau \\
    0 & \text{nếu } \hat{p} < \tau
  \end{cases}
  \label{eq:decision}
\end{equation}
 
Ngưỡng $\tau = 0{,}5$ là lựa chọn mặc định và hợp lý trong điều kiện dữ liệu cân bằng. Tuy nhiên, trong thực tế, khi bài toán yêu cầu độ nhạy cao - ví dụ: phát hiện ung thư, trong đó bỏ sót dương tính thật có chi phí rất cao - ngưỡng $\tau$ có thể được điều chỉnh xuống thấp hơn để tăng độ nhạy (\textit{recall}) nhưng chấp nhận giảm độ chính xác (\textit{precision}). Việc lựa chọn ngưỡng phù hợp thường được hỗ trợ bởi đường cong ROC (\textit{Receiver Operating Characteristic}) và chỉ số AUC (\textit{Area Under Curve}).
 
% ============================================================
\subsubsection{Hàm mất mát và phương pháp tối ưu}
\label{sec:gradient}
% ============================================================
 
Để huấn luyện mô hình, cần tìm bộ tham số $(\mathbf{w},\,b)$ sao cho xác suất dự đoán khớp tốt nhất với nhãn thực tế trong tập dữ liệu. Hồi quy logistic sử dụng hàm mất mát Log-Loss, được dẫn xuất từ nguyên lý ước lượng hợp lý cực đại (\textit{Maximum Likelihood Estimation} --- MLE):
 
\begin{equation}
  \mathcal{L}(\mathbf{w}, b)
  = -\frac{1}{m} \sum_{i=1}^{m}
    \Bigl[
      y^{(i)} \log \hat{p}^{(i)}
      + \bigl(1 - y^{(i)}\bigr) \log \bigl(1 - \hat{p}^{(i)}\bigr)
    \Bigr]
  \label{eq:loss}
\end{equation}
 
\noindent trong đó $m$ là số lượng mẫu huấn luyện, $y^{(i)}$ là nhãn thực tế và $\hat{p}^{(i)}$ là xác suất dự đoán của mẫu thứ $i$.
 
Hàm mất mát này có tính chất lồi (\textit{convex}) đối với không gian tham số, đảm bảo rằng mọi nghiệm cục bộ cũng là nghiệm toàn cục - một tính chất rất được ưa chuộng trong tối ưu hóa. Để cực tiểu hóa $\mathcal{L}$, phương pháp phổ biến nhất là \textit{gradient descent}, các biến thể như \textit{stochastic gradient descent}. Quy tắc cập nhật tham số tại mỗi bước lặp là:
 
\begin{equation}
  \mathbf{w} \leftarrow \mathbf{w} - \eta \cdot \frac{\partial \mathcal{L}}{\partial \mathbf{w}},
  \qquad
  b \leftarrow b - \eta \cdot \frac{\partial \mathcal{L}}{\partial b}
  \label{eq:update}
\end{equation}
 
\noindent trong đó $\eta > 0$ là tốc độ học - một siêu tham số kiểm soát bước dịch chuyển trong không gian tham số.\\
 
Ngoài ra, để tránh hiện tượng mô hình học mọi thứ quá khớp (\textit{overfitting}), người ta thường bổ sung thêm các thành phần chính quy hóa (\textit{regularization}) vào hàm mất mát. Phổ biến nhất là chính quy hóa L2 (\textit{Ridge}):
 
\begin{equation}
  \mathcal{L}_{\mathrm{reg}}
  = \mathcal{L} + \lambda \|\mathbf{w}\|_{2}^{2}
  \label{eq:l2}
\end{equation}
 
\noindent hoặc chính quy hóa L1 (\textit{Lasso}) nhằm khuyến khích thưa thớt trong vector trọng số.
 
% ============================================================
\subsubsection{Ưu điểm}
% ============================================================
 
Hồi quy logistic sở hữu nhiều ưu điểm thiết thực khiến nó vẫn là một trong những mô hình được lựa chọn hàng đầu trong nhiều ứng dụng thực tế. Trước hết, mô hình có tính diễn giải cao: các hệ số $w_i$ phản ánh trực tiếp tác động của từng đặc trưng đến log-odds của kết quả, cho phép chuyên gia lĩnh vực hiểu và kiểm chứng dễ dàng. Đây là ưu điểm quan trọng trong các lĩnh vực đòi hỏi tính minh bạch như y tế, tài chính hay pháp lý.\\
 
Thứ hai, mô hình có độ phức tạp tính toán thấp, phù hợp với các tập dữ liệu lớn và yêu cầu thời gian thực. Hơn nữa, do hàm mất mát có tính lồi, quá trình huấn luyện hội tụ ổn định và không phụ thuộc vào giá trị khởi tạo tham số. Ngoài ra, đầu ra xác suất của mô hình cung cấp thông tin phong phú hơn so với nhãn nhị phân đơn thuần, giúp người dùng linh hoạt hơn trong việc điều chỉnh ngưỡng ra quyết định tùy theo bài toán cụ thể.\\
 
% ============================================================
\subsubsection{Hạn chế}
% ============================================================
 
Bên cạnh các ưu điểm nổi bật, hồi quy logistic cũng tồn tại một số hạn chế đáng kể cần được nhận thức rõ trước khi triển khai. Hạn chế cơ bản nhất là giả định về tính tuyến tính: mô hình chỉ có khả năng học các ranh giới quyết định tuyến tính trong không gian đặc trưng. Do đó, với các bài toán có cấu trúc phi tuyến phức tạp, hồi quy logistic thường cho kết quả kém hơn so với các mô hình phức tạp hơn như mô hình \ref{sec:svm} (với nhân phi tuyến), cây quyết định hay mạng nơ-ron nhân tạo.\\
 
Một hạn chế khác là mô hình nhạy cảm với đa cộng tuyến - hiện tượng các biến đầu vào có tương quan cao với nhau - điều này có thể làm cho ước lượng tham số trở nên không ổn định. Ngoài ra, hồi quy logistic yêu cầu dữ liệu không có nhiễu quá lớn và không xử lý tốt các biến đầu vào có phân phối rất lệch. Cuối cùng, trong bối cảnh dữ liệu mất cân bằng lớp, hiệu suất mô hình có thể bị suy giảm đáng kể nếu không áp dụng các kỹ thuật bổ trợ như \textit{oversampling}, \textit{undersampling} hay điều chỉnh trọng số lớp.\\
 
% ============================================================
\subsubsection{Ứng dụng thực tế}
% ============================================================

Hồi quy logistic phát huy hiệu quả tối ưu trong một số điều kiện và ngữ cảnh cụ thể. Mô hình phù hợp nhất khi biến mục tiêu là nhị phân và quan hệ giữa các đặc trưng đầu vào với log-odds của kết quả có tính xấp xỉ tuyến tính. Mô hình cũng hoạt động tốt khi dữ liệu huấn luyện tương đối cân bằng giữa các lớp, kích thước tập dữ liệu đủ lớn so với số chiều đặc trưng, và các biến đầu vào đã được tiền xử lý kỹ lưỡng (chuẩn hóa, xử lý giá trị thiếu, loại bỏ đa cộng tuyến). Đặc biệt, trong các tình huống đòi hỏi mô hình có khả năng giải thích được - chẳng hạn khi cần trình bày kết quả cho các bên liên quan phi kỹ thuật hoặc tuân thủ các quy định về minh bạch thuật toán - hồi quy logistic là lựa chọn ưu tiên so với các mô hình ``hộp đen'' phức tạp hơn.\\

Nhờ vào sự đơn giản, tốc độ và khả năng diễn giải, hồi quy logistic được ứng dụng rộng rãi trong nhiều lĩnh vực khác nhau. Trong y tế, mô hình được dùng để phân tích nguy cơ mắc bệnh dựa trên các chỉ số lâm sàng, ví dụ dự đoán khả năng tái phát ung thư hoặc chẩn đoán bệnh tim mạch. Trong ngành tài chính và ngân hàng, hồi quy logistic là công cụ kinh điển trong đánh giá tín dụng (\textit{credit scoring}) và phát hiện gian lận giao dịch (\textit{fraud detection}). Trong lĩnh vực tiếp thị kỹ thuật số, mô hình được sử dụng để dự đoán khả năng khách hàng nhấp vào quảng cáo (\textit{click-through rate prediction}) hoặc rời bỏ dịch vụ (\textit{churn prediction}). Trong xử lý ngôn ngữ tự nhiên, hồi quy logistic thường được dùng như một bộ phân loại cơ sở trong các tác vụ phân tích cảm xúc (\textit{sentiment analysis}) hay lọc thư rác (\textit{spam detection}). \\

\subsection{Mô hình SVM (Support Vector Machine)}
% ============================================================
\subsubsection{Khái niệm}
% ============================================================

Support Vector Machine (SVM) là một thuật toán học máy có giám sát được phát triển bởi Vladimir Vapnik và Alexey Chervonenkis vào những năm 1960, sau đó được hoàn thiện và phổ biến rộng rãi thông qua công trình của Cortes và Vapnik năm 1995. SVM được xây dựng trên nền tảng lý thuyết học thống kê, cụ thể là lý thuyết về tối thiểu hóa rủi ro có cấu trúc (\textit{structural risk minimization - SRM}), nhằm tìm kiếm một mô hình không chỉ khớp tốt với dữ liệu huấn luyện mà còn có khả năng tổng quát hóa cao trên dữ liệu chưa thấy.\\

Về bản chất, SVM là một bộ phân loại tuyến tính tìm kiếm siêu phẳng tối ưu phân tách các lớp dữ liệu trong không gian đặc trưng, sao cho khoảng cách từ siêu phẳng đến các điểm dữ liệu gần nhất của mỗi lớp là lớn nhất có thể. Thông qua kỹ thuật hàm nhân, SVM có thể mở rộng sang các bài toán phi tuyến mà không cần tường minh tính toán trong không gian chiều cao.\\

Cũng như hồi quy Logistic, SVM được thiết kế chủ yếu cho các bài toán phân loại nhị phân, trong đó mục tiêu là phân chia không gian đặc trưng thành hai vùng tương ứng với hai nhãn lớp. Ngoài ra, thông qua các phương pháp mở rộng như \textit{One-vs-Rest} (OvR) và \textit{One-vs-One} (OvO), SVM có thể xử lý hiệu quả các bài toán phân loại đa lớp. Thêm vào đó, biến thể \textit{Support Vector Regression} (SVR) cho phép áp dụng nguyên lý của SVM vào các bài toán hồi quy, mở rộng đáng kể phạm vi ứng dụng của mô hình.\\

Điểm nổi bật của SVM so với nhiều thuật toán phân loại khác nằm ở khả năng hoạt động hiệu quả trong không gian chiều cao và trên các tập dữ liệu có kích thước vừa và nhỏ - đặc biệt khi số chiều đặc trưng vượt quá số lượng mẫu, một tình huống mà nhiều thuật toán khác gặp khó khăn.\\

% ============================================================
\subsubsection{Siêu Phẳng, Margin và Support Vectors}
% ============================================================

Trong không gian $n$ chiều $\mathbb{R}^n$, một siêu phẳng là tập hợp các điểm $\mathbf{x}$ thỏa mãn phương trình tuyến tính:

\begin{equation}
  \mathbf{w}^T \mathbf{x} + b = 0
  \label{eq:hyperplane}
\end{equation}

\noindent trong đó $\mathbf{w} \in \mathbb{R}^n$ là vector pháp tuyến xác định hướng của siêu phẳng và $b \in \mathbb{R}$ là hệ số chệch xác định vị trí của siêu phẳng trong không gian. Siêu phẳng này chia không gian thành hai nửa tương ứng với hai lớp phân loại: lớp dương tính ($y = +1$) khi $\mathbf{w}^T \mathbf{x} + b > 0$ và lớp âm tính ($y = -1$) khi $\mathbf{w}^T \mathbf{x} + b < 0$.

Khoảng cách từ một điểm $\mathbf{x}_i$ đến siêu phẳng được xác định bởi:

\begin{equation}
  d_i = \frac{|\mathbf{w}^T \mathbf{x}_i + b|}{\|\mathbf{w}\|}
  \label{eq:distance}
\end{equation}

Biên độ (\textit{margin}) của siêu phẳng được định nghĩa là tổng khoảng cách từ siêu phẳng đến điểm gần nhất của mỗi lớp. Đối với siêu phẳng được chuẩn hóa sao cho $|\mathbf{w}^T \mathbf{x}_i + b| = 1$ với các điểm gần nhất, biên độ tổng bằng $\dfrac{2}{\|\mathbf{w}\|}$.

Các \textit{support vectors} là những điểm dữ liệu nằm trên biên giới của biên độ, tức là những điểm thỏa mãn $|\mathbf{w}^T \mathbf{x}_i + b| = 1$. Đây là các điểm dữ liệu duy nhất quyết định vị trí và hướng của siêu phẳng tối ưu --- nếu loại bỏ các điểm này, siêu phẳng sẽ thay đổi, trong khi loại bỏ các điểm còn lại thì không.

% ============================================================
\subsubsection{Tối Đa Hóa Khoảng Cách Biên}
% ============================================================

Mục tiêu của SVM là tìm siêu phẳng tối đa hóa biên độ $\dfrac{2}{\|\mathbf{w}\|}$, tương đương với việc tối thiểu hóa $\|\mathbf{w}\|^2$, đồng thời đảm bảo tất cả các điểm dữ liệu được phân loại đúng. Lập luận cơ bản đằng sau nguyên tắc này là: trong số tất cả các siêu phẳng có thể phân tách hai lớp, siêu phẳng có biên độ lớn nhất cung cấp sự phân tách ``an toàn'' nhất, làm giảm thiểu rủi ro phân loại sai trên dữ liệu mới. Đây chính là hiện thân của nguyên tắc tối thiểu hóa rủi ro có cấu trúc mà SVM được xây dựng dựa trên.


\subsubsection{Nguyên lý hoạt động}
% ============================================================

Ý tưởng cốt lõi của SVM là tìm một ranh giới quyết định (\textit{decision boundary}) phân tách hai lớp dữ liệu sao cho khoảng cách từ ranh giới đó đến các điểm dữ liệu gần nhất của mỗi lớp được tối đa hóa. Nguyên lý này xuất phát từ lập luận rằng một ranh giới quyết định có biên độ lớn hơn sẽ có độ tin cậy cao hơn và ít nhạy cảm hơn với nhiễu trong dữ liệu, từ đó dẫn đến khả năng tổng quát hóa tốt hơn trên dữ liệu mới.\\

Quá trình huấn luyện của SVM về bản chất là một bài toán tối ưu hóa lồi, trong đó mô hình tìm kiếm các bộ tham số xác định siêu phẳng tối ưu thông qua phương pháp \textit{gradient descent} - phương pháp lặp cập nhật tham số theo chiều giảm nhanh nhất của hàm mục tiêu. Điều quan trọng cần nhấn mạnh là SVM chỉ phụ thuộc vào một tập con nhỏ các điểm dữ liệu huấn luyện - gọi là các vector hỗ trợ (\textit{support vectors}) --- để xác định ranh giới quyết định, thay vì toàn bộ tập dữ liệu. Tính chất này mang lại cho SVM sự hiệu quả đặc biệt về mặt bộ nhớ và độ bền trước các điểm dữ liệu ngoại lệ nằm xa ranh giới.


\textbf{Hard-Margin SVM}

Trong trường hợp dữ liệu tuyến tính phân tách được hoàn toàn (\textit{linearly separable}), bài toán tối ưu hóa của SVM được phát biểu như sau:

\begin{equation}
  \min_{\mathbf{w},\, b} \quad \frac{1}{2}\|\mathbf{w}\|^2
  \label{eq:hardmargin_obj}
\end{equation}

\begin{equation}
  \text{subject to} \quad y_i\!\left(\mathbf{w}^T \mathbf{x}_i + b\right) \geq 1,
  \quad \forall\, i = 1, \ldots, m
  \label{eq:hardmargin_const}
\end{equation}

\noindent trong đó $y_i \in \{-1, +1\}$ là nhãn lớp của mẫu thứ $i$, $\mathbf{x}_i \in \mathbb{R}^n$ là vector đặc trưng, $\mathbf{w}$ là vector trọng số và $b$ là hệ số chệch. Ràng buộc \eqref{eq:hardmargin_const} đảm bảo rằng mọi điểm dữ liệu nằm đúng phía và cách xa siêu phẳng ít nhất một khoảng bằng $\dfrac{1}{\|\mathbf{w}\|}$.

\textbf{Soft-Margin SVM}

Trong thực tế, dữ liệu hiếm khi phân tách tuyến tính hoàn hảo. Để xử lý nhiễu và các điểm ngoại lệ, SVM được mở rộng thành biến thể \textit{soft-margin} bằng cách đưa vào các biến bù (\textit{slack variables}) $\xi_i \geq 0$:

\begin{equation}
  \min_{\mathbf{w},\, b,\, \boldsymbol{\xi}} \quad
  \frac{1}{2}\|\mathbf{w}\|^2 + C \sum_{i=1}^{m} \xi_i
  \label{eq:softmargin_obj}
\end{equation}

\begin{equation}
  \text{subject to} \quad y_i\!\left(\mathbf{w}^T \mathbf{x}_i + b\right)
  \geq 1 - \xi_i, \quad \xi_i \geq 0, \quad \forall\, i
  \label{eq:softmargin_const}
\end{equation}

\noindent Tham số $C > 0$ là siêu tham số điều tiết (\textit{regularization parameter}) kiểm soát sự đánh đổi giữa việc tối đa hóa biên độ và chấp nhận vi phạm phân loại. Khi $C$ lớn, mô hình ưu tiên phân loại đúng tất cả các điểm (ít vi phạm hơn, biên độ nhỏ hơn); khi $C$ nhỏ, mô hình ưu tiên biên độ lớn và chấp nhận một số vi phạm.

\subsubsubsection{Bài Toán Đối Ngẫu}

Bài toán tối ưu trên thường được giải thông qua dạng đối ngẫu Lagrange (\textit{Lagrangian dual}), dẫn đến bài toán:

\begin{equation}
  \max_{\boldsymbol{\alpha}} \quad
  \sum_{i=1}^{m} \alpha_i
  - \frac{1}{2} \sum_{i=1}^{m} \sum_{j=1}^{m}
    \alpha_i \alpha_j y_i y_j\, \mathbf{x}_i^T \mathbf{x}_j
  \label{eq:dual_obj}
\end{equation}

\begin{equation}
  \text{subject to} \quad 0 \leq \alpha_i \leq C,
  \quad \sum_{i=1}^{m} \alpha_i y_i = 0
  \label{eq:dual_const}
\end{equation}

\noindent trong đó $\alpha_i \geq 0$ là các nhân tử Lagrange. Nghiệm của bài toán đối ngẫu cho phép biểu diễn vector trọng số tối ưu như tổ hợp tuyến tính của các support vectors:

\begin{equation}
  \mathbf{w}^* = \sum_{i=1}^{m} \alpha_i y_i \mathbf{x}_i
  \label{eq:w_star}
\end{equation}

\noindent Quan trọng hơn, dạng đối ngẫu cho thấy tích vô hướng $\mathbf{x}_i^T \mathbf{x}_j$ có thể được thay thế bằng hàm nhân $K(\mathbf{x}_i, \mathbf{x}_j)$, mở đường cho kỹ thuật kernel.

% ============================================================
\subsubsection{Linear SVM và Non-linear SVM}
% ============================================================

\textbf{Linear SVM} áp dụng trực tiếp công thức trên trong không gian đặc trưng gốc, phù hợp khi dữ liệu có thể phân tách tuyến tính hoặc xấp xỉ tuyến tính. Mô hình này đơn giản, nhanh và dễ diễn giải, nhưng bị giới hạn khi ranh giới quyết định thực sự có dạng phi tuyến.

\textbf{Non-linear SVM} giải quyết hạn chế này bằng cách ánh xạ dữ liệu sang một không gian đặc trưng chiều cao hơn $\mathcal{H}$ thông qua một hàm ánh xạ $\phi: \mathbb{R}^n \to \mathcal{H}$, trong đó các lớp có thể phân tách tuyến tính. Kỹ thuật kernel cho phép thực hiện điều này một cách ngầm định mà không cần tường minh tính toán $\phi(\mathbf{x})$, tránh được gánh nặng tính toán trong không gian chiều rất cao hoặc vô hạn.

% ============================================================
\subsubsection{Kernel trick và các hàm kernel phổ biến}
% ============================================================

Kỹ thuật kernel (\textit{kernel trick}) dựa trên quan sát rằng trong bài toán đối ngẫu \eqref{eq:dual_obj}, dữ liệu chỉ xuất hiện dưới dạng tích vô hướng $\mathbf{x}_i^T \mathbf{x}_j$. Bằng cách thay thế tích này bằng hàm nhân $K(\mathbf{x}_i, \mathbf{x}_j) = \phi(\mathbf{x}_i)^T \phi(\mathbf{x}_j)$, mô hình ngầm định hoạt động trong không gian chiều cao mà không cần tường minh tính toán $\phi$.

\textbf{Linear Kernel} là trường hợp đơn giản nhất, tương ứng với SVM tuyến tính thông thường:
\begin{equation}
  K(\mathbf{x}_i, \mathbf{x}_j) = \mathbf{x}_i^T \mathbf{x}_j
  \label{eq:linear_kernel}
\end{equation}

\textbf{Polynomial Kernel} ánh xạ dữ liệu vào không gian đa thức bậc $d$, cho phép mô hình học các ranh giới quyết định có dạng đa thức:
\begin{equation}
  K(\mathbf{x}_i, \mathbf{x}_j)
  = \left(\gamma\, \mathbf{x}_i^T \mathbf{x}_j + r\right)^d,
  \quad \gamma > 0,\; d \in \mathbb{Z}^+
  \label{eq:poly_kernel}
\end{equation}

\textbf{Radial Basis Function Kernel} (RBF, hay Gaussian Kernel) là hàm nhân được sử dụng phổ biến nhất trong thực tế, ánh xạ dữ liệu vào không gian vô hạn chiều:
\begin{equation}
  K(\mathbf{x}_i, \mathbf{x}_j)
  = \exp\!\left(-\gamma \|\mathbf{x}_i - \mathbf{x}_j\|^2\right),
  \quad \gamma > 0
  \label{eq:rbf_kernel}
\end{equation}

\noindent Tham số $\gamma$ kiểm soát phạm vi ảnh hưởng của mỗi điểm dữ liệu: $\gamma$ lớn dẫn đến mô hình phức tạp hơn (ranh giới quyết định cục bộ hơn), trong khi $\gamma$ nhỏ tạo ra ranh giới mượt mà hơn.

\textbf{Sigmoid Kernel} lấy cảm hứng từ hàm kích hoạt trong mạng nơ-ron:
\begin{equation}
  K(\mathbf{x}_i, \mathbf{x}_j)
  = \tanh\!\left(\gamma\, \mathbf{x}_i^T \mathbf{x}_j + r\right)
  \label{eq:sigmoid_kernel}
\end{equation}

\noindent Kernel này không phải lúc nào cũng thỏa mãn điều kiện Mercer --- điều kiện để đảm bảo bài toán tối ưu lồi --- do đó ít được sử dụng hơn trong thực tế so với RBF hay Polynomial.

% ============================================================
\subsubsection{Hàm mất mát và phương pháp tối ưu hóa}
% ============================================================

SVM sử dụng hàm mất mát \textit{Hinge Loss}, được định nghĩa cho từng mẫu như sau:

\begin{equation}
  \ell_i = \max\!\left(0,\; 1 - y_i\!\left(\mathbf{w}^T \mathbf{x}_i + b\right)\right)
  \label{eq:hinge}
\end{equation}

\noindent Hàm mất mát này bằng $0$ khi điểm dữ liệu được phân loại đúng và nằm ngoài vùng biên độ; ngược lại, nó tăng tuyến tính theo mức độ vi phạm. Hàm mục tiêu tổng thể của soft-margin SVM có thể được viết lại dưới dạng:

\begin{equation}
  \mathcal{L}(\mathbf{w}, b)
  = \frac{1}{2}\|\mathbf{w}\|^2
  + C \sum_{i=1}^{m} \max\!\left(0,\;
    1 - y_i\!\left(\mathbf{w}^T \mathbf{x}_i + b\right)\right)
  \label{eq:total_loss}
\end{equation}

\noindent Đây là hàm lồi nhưng không khả vi tại điểm $1 - y_i(\mathbf{w}^T \mathbf{x}_i + b) = 0$, do đó không thể áp dụng trực tiếp gradient descent thông thường. Phương pháp tối ưu phổ biến nhất cho SVM là \textit{Sequential Minimal Optimization} (SMO) --- một thuật toán chuyên biệt phân rã bài toán lập trình toàn phương (\textit{Quadratic Programming} --- QP) lớn thành các bài toán QP con hai biến và giải chúng theo phương pháp giải tích. Ngoài ra, \textit{subgradient descent} và các phương pháp tối ưu dựa trên nhân tử Lagrange cũng được sử dụng rộng rãi.

% ============================================================
\subsubsection{Ưu điểm và hạn chế}
% ============================================================

SVM mang lại nhiều ưu thế đáng kể trong các bài toán phân loại. Trước hết, mô hình có khả năng tổng quát hóa tốt nhờ nguyên tắc tối đa hóa biên độ, giúp giảm thiểu rủi ro quá khớp (\textit{overfitting}) ngay cả trong không gian chiều cao. Thứ hai, thông qua kỹ thuật kernel, SVM có thể xử lý hiệu quả các bài toán phi tuyến phức tạp mà không bị ảnh hưởng bởi lời nguyền chiều cao (\textit{curse of dimensionality}) theo cách tương tự như các phương pháp khác. Thứ ba, nghiệm của SVM là duy nhất và toàn cục do bài toán tối ưu có tính lồi, loại bỏ rủi ro hội tụ về cực tiểu cục bộ. Cuối cùng, do chỉ phụ thuộc vào support vectors, mô hình hiệu quả về bộ nhớ và bền vững trước các điểm dữ liệu không ảnh hưởng đến ranh giới quyết định.\\

Mặc dù có nhiều ưu điểm nổi bật, SVM cũng tồn tại những hạn chế cần cân nhắc kỹ trước khi áp dụng. Hạn chế lớn nhất là chi phí tính toán cao khi huấn luyện trên tập dữ liệu lớn: độ phức tạp thời gian của các thuật toán giải QP thông thường là $O(m^2)$ đến $O(m^3)$ theo số mẫu $m$, khiến SVM trở nên kém hiệu quả khi $m$ lên đến hàng triệu. Thứ hai, việc lựa chọn kernel và điều chỉnh các siêu tham số ($C$, $\gamma$, $d$) đòi hỏi kinh nghiệm và thường phải dựa vào tìm kiếm lưới (\textit{grid search}) kết hợp kiểm định chéo (\textit{cross-validation}), làm tăng chi phí tính toán tổng thể. Thứ ba, không giống như hồi quy logistic, SVM không trực tiếp cung cấp ước lượng xác suất --- xác suất lớp có thể được tính toán thông qua phương pháp Platt scaling, nhưng điều này đòi hỏi thêm bước hiệu chỉnh sau huấn luyện. Cuối cùng, kết quả của SVM khó diễn giải hơn so với các mô hình tuyến tính đơn giản, đặc biệt khi sử dụng kernel phi tuyến.

% ============================================================
\subsubsection{Ứng dụng thực tế}
% ============================================================

SVM đã được ứng dụng thành công trong nhiều lĩnh vực đòi hỏi độ chính xác cao. Trong nhận dạng ảnh và thị giác máy tính, SVM từng là phương pháp hàng đầu trong phân loại ảnh và nhận dạng chữ viết tay trước khi học sâu (\textit{deep learning}) trở nên phổ biến. Trong phân loại văn bản và xử lý ngôn ngữ tự nhiên, SVM với kernel tuyến tính cho kết quả xuất sắc trong phân loại thư rác và phân tích cảm xúc do tính hiệu quả trong không gian chiều cao thưa (\textit{sparse high-dimensional space}). Trong y sinh học, SVM được ứng dụng rộng rãi trong phân loại dữ liệu gen, chẩn đoán bệnh từ dữ liệu y tế, và phân tích tín hiệu não (EEG/fMRI). Trong tài chính, mô hình được sử dụng cho phát hiện gian lận và dự đoán biến động thị trường.\\

SVM phát huy hiệu quả cao nhất trong một số điều kiện đặc thù. Mô hình phù hợp khi tập dữ liệu có kích thước vừa và nhỏ (dưới vài chục nghìn mẫu), số chiều đặc trưng lớn hoặc thậm chí vượt quá số mẫu, và ranh giới quyết định có tính phi tuyến rõ ràng. SVM cũng là lựa chọn tốt khi dữ liệu có biên độ phân tách rõ ràng giữa các lớp và khi yêu cầu một nghiệm tối ưu toàn cục được đảm bảo. Ngược lại, SVM không phải lựa chọn lý tưởng cho các tập dữ liệu rất lớn, các bài toán yêu cầu ước lượng xác suất trực tiếp, hoặc khi thời gian huấn luyện là yếu tố quan trọng.